\documentclass[11pt]{article}
\usepackage[T1]{fontenc}
\usepackage{lmodern,amsmath,amssymb,amsthm}
\usepackage[margin=1in]{geometry}
\usepackage[hidelinks]{hyperref}
\newtheorem{theorem}{Theorem}
\newtheorem{lemma}{Lemma}
\newtheorem{corollary}{Corollary}
\newcommand{\GV}{\operatorname{gv}}
\newcommand{\Diff}{\operatorname{Diff}}
\title{Unbounded Godbillon--Vey numbers on a fixed hyperbolic manifold}
\author{Alec Kriebel\\[2pt]
\small Independent researcher\\
\small\url{https://orcid.org/0009-0001-9320-500X}}
\date{13 October 2026}
\begin{document}
\maketitle
\begin{abstract}
Classical compact-support suspension theory yields a smooth foliated
surface-times-interval with product endpoint collars and nonzero
Godbillon--Vey integral. We place arbitrarily many copies inside one fixed
pseudo-Anosov mapping torus. One embedded closed transversal meets every
leaf, and the Godbillon--Vey numbers grow linearly with the number of
copies. This gives a negative answer to the unrestricted volume-bound
question printed as Question~10.6 in Calegari's 2002 manuscript.
The construction uses classical ingredients; historical priority of this
specific consequence remains unestablished.
\end{abstract}

\section{Statement and scope}
For a smooth cooriented codimension-one foliation $\mathcal F$ of a closed
oriented three-manifold, choose a defining form $\alpha$ and a form $\eta$
with $d\alpha=\alpha\wedge\eta$. We use the ordinary Godbillon--Vey number
\[
       \GV(\mathcal F)=\int_M\eta\wedge d\eta .
\]
The closed integral is independent of these choices. We call $\mathcal F$
taut when a closed transversal meets every leaf.

\begin{theorem}\label{thm:main}
There exist a closed connected oriented hyperbolic three-manifold $M$,
a number $a>0$, and an embedded smooth circle $\gamma\subset M$ such that,
for every integer $N\geq1$, $M$ admits a smooth cooriented taut foliation
$\mathcal F_N$ with
\[
             \GV(\mathcal F_N)=Na .
\]
The same $\gamma$ is positively transverse to, and meets every leaf of,
each $\mathcal F_N$.
\end{theorem}

Fixing one hyperbolic metric on $M$ fixes its positive finite volume.
The theorem therefore rules out any finite bound determined solely by
that volume, including a bound on $|\GV|$. It answers the literal
Question~10.6 of \cite[p.~25]{C02}, using the single-circle definition
in \cite[Definition~1.1]{C02}. Compact leaves are permitted.
This note makes no assertion for minimal foliations, a condition imposed
in the different Question~13.1 of that source.

\paragraph{Attribution and priority.}
The nonzero suspension input is classical Mather--Thurston theory,
with the precise compatibility recorded by Tsuboi~\cite{T81}.
Interval-bundle replacement is classical; see
Gabai~\cite[Operation~2.1.3]{G92}. Hyperbolization is due to
Thurston~\cite{Th98}. A bounded primary-source audit did not authenticate
an earlier explicit statement of Theorem~\ref{thm:main}, but does not
establish its novelty, firstness, or present-day open status.
Several relevant historical full texts remained inaccessible or unread.
Our target is the specified 2002 manuscript version; the final 2003
chapter was not compared. Detailed read-depths and gaps accompany this
note.

\section{A classical suspension block}
Write $I=[0,1]$ and let $G=\Diff^\infty_c(\mathbb R)$, regarded as a
discrete group. These diffeomorphisms preserve orientation.

\begin{lemma}\label{lem:block}
There are a closed connected oriented surface $\Sigma$ of genus at least
two and a smooth cooriented foliation $\mathcal H$ on $\Sigma\times I$,
transverse to every interval fiber and product on endpoint collars,
with defining forms satisfying
\[
 \alpha(\partial_s)=1,\qquad d\alpha=\alpha\wedge\eta,\qquad
 (\alpha,\eta)=(ds,0)\ \text{on both collars},\qquad
 \int_{\Sigma\times I}\eta\wedge d\eta=a>0 .
\]
Every leaf of $\mathcal H$ meets the interval fiber over each $p\in\Sigma$.
\end{lemma}

\begin{proof}
We first specify the classical input. Tsuboi~\cite[pp.~1--3,6]{T81}
records the surjective smooth GV map
$H_3(B\overline\Gamma^\infty_1;\mathbb Z)\to\mathbb R$, Mather's
isomorphism from $H_2(G;\mathbb Z)$ to this homology group, and the
commutative diagram identifying that map with compactification into
circle diffeomorphisms followed by the actual circle-bundle suspension.
Integral group-homology two-classes are represented by closed oriented
surfaces. Choose a class with positive GV suspension value, then a
positive connected component of a surface representative. Adding
trivially mapped handles, via a degree-one collapse, preserves its
class and arranges genus at least two. Thus one representation
$\rho:\pi_1\Sigma\to G$ has compactified circle suspension of value $a>0$.
No uniform genus bound over all GV values is needed.

Finitely many generator images have supports in a common compact
interval. Every word fixes its complement. Compactifying $\mathbb R$
to $(0,1)$ gives a suspension
\[
         E=(\widetilde\Sigma\times I)/\pi_1\Sigma
\]
with common product collars. It admits a smooth fiber trivialization
relative to those collars: take increasing local fiber coordinates
$s_i$ agreeing with the collar coordinate and a base partition of unity
$\lambda_i$. The average $\sum_i\lambda_i s_i$ has positive vertical
derivative, endpoint values $0,1$, and the prescribed collar values.
It is consequently a smooth coordinate on every fiber, with smooth
inverse. This trivializes the bundle, while preserving its possibly
nontrivial horizontal foliation.

In this trivialization let $V=\partial_s$ and normalize the cooriented
defining form by $\alpha(V)=1$. Then $\alpha=ds$ on the collars.
Frobenius and contraction give
\[
  0=\iota_V(\alpha\wedge d\alpha)
    =d\alpha-\alpha\wedge\iota_Vd\alpha .
\]
Set $\eta=\iota_Vd\alpha$; it vanishes on the collars.
Complete $I$ to a circle with an additional product interval, extending
$\alpha$ by the circle-coordinate form and $\eta$ by zero.
This recovers the classical compactified suspension. The added region
has zero GV integrand, proving the asserted block integral.

Finally the suspension leaf through $u\in I$ is
$\widetilde\Sigma/\operatorname{Stab}(u)$ and projects as a covering
onto $\Sigma$. Hence every leaf meets every prescribed interval fiber.
\end{proof}

\section{Fixed ambient manifold and finite stacking}
\begin{proof}[Proof of Theorem~\ref{thm:main}]
Choose an orientation-preserving pseudo-Anosov mapping class on $\Sigma$.
A smooth representative $\varphi$ can fix a small disk pointwise.
For example, the filling-multicurve twist construction
\cite[Theorem~3.1]{P88}, also proved in \cite[pp.~467--468]{F92},
uses smooth twists in thin annuli leaving a complementary disk fixed.
The pseudo-Anosov property is required of the mapping class; the chosen
disk-fixed representative need not itself be a pseudo-Anosov normal form.
Fix $\varphi$ and
\[
 M=(\Sigma\times[0,1])/\bigl((x,1)\sim(\varphi(x),0)\bigr).
\]
Mapping-torus hyperbolization \cite[Theorem~0.1, Proposition~2.6,
\S5]{Th98} makes this closed oriented manifold hyperbolic.
Fix one such metric once and for all.

Take $p$ in the fixed disk. The curve $\gamma(t)=[p,t]$ is a smooth
embedded circle, and $V=\partial_t$ descends through the seam.
For a given finite $N$, choose pairwise disjoint closed intervals
$[a_j,b_j]\subset(0,1)$ with positive intervening gaps, and write
$\ell_j=b_j-a_j>0$. On the $j$th slab use
\[
 h_j(x,t)=\left(x,\frac{t-a_j}{\ell_j}\right),\qquad
 \alpha_j=\ell_j h_j^*\alpha,\qquad \eta_j=h_j^*\eta .
\]
They satisfy $d\alpha_j=\alpha_j\wedge\eta_j$ and
$\alpha_j(V)=1$. On open endpoint collars they equal $dt,0$,
with equality of all derivatives. Extend by $dt,0$ on the complement,
including the mapping-torus seam. The resulting smooth nonvanishing
form $\alpha_N$ is integrable and defines a cooriented foliation
$\mathcal F_N$. The manifold, metric, and seam identification have not
changed.

In each slab every suspension leaf meets the vertical interval over $p$,
by Lemma~\ref{lem:block}. In gaps, collars, and the seam neighborhood
every leaf is a surface fiber and meets $\gamma$. Interior block leaves
cannot cross a product collar leaf, so this accounts for every leaf.
Also $\alpha_N(V)=1$ everywhere. Thus the same $\gamma$ is a positive
closed transversal meeting every leaf, proving tautness.
No condition intertwining $\rho$ with $\varphi$ is needed: insertion
takes place in interior slabs with product behavior near the seam.

The form $\eta_N$ vanishes off the inserted blocks and near all
interfaces. Each $h_j$ preserves orientation, so
\[
 \int_{\Sigma\times[a_j,b_j]}\eta_j\wedge d\eta_j
       =\int_{\Sigma\times I}\eta\wedge d\eta=a .
\]
Summing gives $\GV(\mathcal F_N)=Na$. Only finite stacks are used;
derivatives may grow with $N$ while each foliation remains smooth.
\end{proof}

\paragraph{Verification and limitations.}
The existence proof uses the cited classical suspension and
hyperbolization theorems. It supplies no explicit numerical holonomy
representation or certified hyperbolic metric.
The accompanying exact-rational and symbolic controls check local
form identities, compression, and finite bookkeeping; they do not
replace the global proof. Their explicit failure checks remain active
under Python optimization. Historical verifier deficiencies, independent
adversarial audits, and the bounded priority search are documented
separately. All primary-source PDFs are excluded from the supplement.

\paragraph{AI and review disclosure.}
AI tools were used extensively in development, drafting, source checking,
computational reproduction, and adversarial review. This preprint is
unrefereed and has not undergone conventional human peer review.

\begin{thebibliography}{9}
\bibitem{C02} D.~Calegari,
\emph{Problems in foliations and laminations of 3-manifolds},
arXiv:math/0209081v1 (2002), Question~10.6.
Published citation: Proc.\ Sympos.\ Pure Math.\ \textbf{71} (2003), 297--335;
\href{https://doi.org/10.1090/pspum/071/2024640}{doi:10.1090/pspum/071/2024640}.
\bibitem{T81} T.~Tsuboi,
\emph{On 2-cycles of $B\Diff(S^1)$ which are represented by foliated
$S^1$-bundles over $T^2$},
Ann.\ Inst.\ Fourier \textbf{31}(2) (1981), 1--59;
\href{https://doi.org/10.5802/aif.828}{doi:10.5802/aif.828}.
\bibitem{G92} D.~Gabai,
\emph{Taut foliations of 3-manifolds and suspensions of $S^1$},
Ann.\ Inst.\ Fourier \textbf{42}(1--2) (1992), 193--208;
\href{https://doi.org/10.5802/aif.1289}{doi:10.5802/aif.1289}.
\bibitem{P88} R.~C.~Penner,
\emph{A construction of pseudo-Anosov homeomorphisms},
Trans.\ Amer.\ Math.\ Soc.\ \textbf{310}(1) (1988), 179--197;
\href{https://doi.org/10.1090/S0002-9947-1988-0930079-9}
{doi:10.1090/S0002-9947-1988-0930079-9}.
\bibitem{F92} A.~Fathi,
\emph{D\'emonstration d'un th\'eor\`eme de Penner sur la composition
des twists de Dehn},
Bull.\ Soc.\ Math.\ France \textbf{120}(4) (1992), 467--484.
\bibitem{Th98} W.~P.~Thurston,
\emph{Hyperbolic structures on 3-manifolds, II: Surface groups and
3-manifolds which fiber over the circle},
August 1986 preprint, January 1998 eprint,
\href{https://arxiv.org/abs/math/9801045v1}{arXiv:math/9801045v1}.
\end{thebibliography}
\end{document}

